\setlength{\baselineskip}{20pt}
\chapter{相关工作}
\label{cha:chap2}

本章系统综述与本文研究密切相关的基础理论与技术方法。首先介绍生成式扩散模型的基本原理与发展历程；随后回顾判别式表征学习中骨干网络的演进脉络；接着综述扩散模型在判别任务中的已有探索；然后讨论高效扩散采样方法；最后介绍知识蒸馏的基本范式与典型应用。

\section{生成式扩散模型}

生成式模型旨在学习数据的底层分布，进而能够合成符合该分布的新样本。典型的生成式模型通过估计每个类别的条件似然$P(\mathbf{x}|y=k)$及先验类别概率$P(y=k)$，再由全概率公式计算边缘分布$P(\mathbf{x})$，从而近似各类别数据的概率分布，并通过从学习到的分布中采样来生成新数据。

经典的生成式模型包括生成对抗网络（Generative Adversarial Networks, GANs）\cite{goodfellow2014generative,karras2019style,mirza2014conditional}和变分自编码器（Variational Autoencoders, VAEs）\cite{kingma2013auto,van2017neural,xu2022socialvae,zhao2017learning}。GAN通过生成器与判别器之间的对抗训练，在隐空间中学习数据的生成过程，能够产生高质量的合成样本。VAE则通过编码器-解码器结构学习数据的隐变量分布，在图像生成和表示学习中均取得了显著进展。这两类方法在建模数据生成过程方面展现出了优秀的性能。

近年来，研究关注逐步转向扩散模型（Diffusion Models）。扩散模型的基本框架最早由Sohl-Dickstein等人\cite{sohl2015deep}于2015年提出，其核心思想是通过迭代去除训练图像中的高斯噪声来建模数据分布。具体而言，扩散模型包含前向扩散和反向去噪两个对称过程。在前向扩散过程中，模型通过马尔可夫链逐步向数据$\mathbf{x}_0$添加高斯噪声，经过$T$步后将原始数据转化为近似的标准高斯分布：
\begin{equation}
q(\mathbf{x}_t|\mathbf{x}_{t-1}) = \mathcal{N}(\mathbf{x}_t; \sqrt{1-\beta_t}\mathbf{x}_{t-1}, \beta_t\mathbf{I})
\end{equation}
其中$\beta_t$为预先设定的噪声调度参数。在反向去噪过程中，模型训练一个参数化的去噪网络$\epsilon_\theta$来逐步恢复原始数据分布：
\begin{equation}
p_\theta(\mathbf{x}_{t-1}|\mathbf{x}_t) = \mathcal{N}(\mathbf{x}_{t-1}; \boldsymbol{\mu}_\theta(\mathbf{x}_t, t), \boldsymbol{\Sigma}_\theta(\mathbf{x}_t, t))
\end{equation}
通过训练去噪网络预测每步添加的噪声$\epsilon$，模型能够从纯噪声出发通过迭代去噪生成高质量的样本。

2020年，Ho等人\cite{ho2020denoising}提出去噪扩散概率模型（DDPM），通过精心设计的噪声调度策略和简化的训练目标函数$\mathcal{L} = \mathbb{E}_{t, \mathbf{x}_0, \epsilon}[||\epsilon - \epsilon_\theta(\mathbf{x}_t, t)||^2]$，使扩散模型在图像生成质量上首次超越GAN，确立了扩散模型作为生成式建模主流方法的地位。Nichol和Dhariwal\cite{nichol2021improved}进一步提出改进型DDPM，通过学习反向过程的方差参数和引入余弦噪声调度策略，在保持样本质量的同时大幅减少了所需的采样步数并提升了对数似然度。Dhariwal和Nichol\cite{dhariwal2021diffusion}则通过系统性的架构搜索和分类器引导（Classifier Guidance）策略，首次在条件图像生成任务上使扩散模型在FID指标上全面超越GAN，奠定了扩散模型在视觉生成领域的主导地位。

在理论框架方面，Song等人\cite{song2021score}提出基于随机微分方程（SDE）的统一生成建模框架，将DDPM和基于分数的生成模型纳入连续时间SDE的统一视角。该工作引入了预测器-校正器框架，并推导出等价的常微分方程（ODE）形式以实现精确的似然计算和高效采样，为扩散模型提供了更加系统化的数学基础。

此后，扩散模型在多种生成任务中持续展现出卓越的能力。Stable Diffusion\cite{rombach2022high}通过将扩散过程迁移至隐空间（Latent Space）进行，实现了高分辨率图像的高效合成。DALL·E\cite{ramesh2021zero}系列将文本编码器与扩散生成器相结合，实现了零样本文本到图像的生成。Imagen\cite{saharia2022photorealistic}通过引入大规模语言模型的深度文本理解能力，进一步提升了文本引导图像生成的逼真度。

在网络架构方面，早期扩散模型普遍采用U-Net作为去噪骨干网络。Peebles和Xie\cite{peebles2023scalable}提出了扩散Transformer（DiT），首次以Transformer架构替代U-Net作为隐空间扩散模型的骨干网络。DiT通过在隐空间图像块上操作，并利用自适应层归一化（AdaLN）注入条件信息，展现出优异的可扩展性——更高的计算量（以GFLOPs衡量）持续带来更低的FID分数。其最大的DiT-XL/2模型在class-conditional ImageNet 256$\times$256基准上取得了2.27的FID，刷新了当时的生成质量记录。这一工作深刻影响了后续大规模视觉生成模型的架构设计。

此外，Lipman等人\cite{lipman2023flow}提出了流匹配（Flow Matching）方法，为训练连续归一化流（CNF）提供了一种无需模拟的高效范式。流匹配通过回归固定条件概率路径的向量场来训练CNF，兼容包括扩散路径在内的多种高斯概率路径族，并发现使用最优传输位移插值定义的路径相比扩散路径具有更高的训练与采样效率。扩散模型还被成功拓展至视频生成\cite{ho2022imagen}、分子结构生成\cite{hoogeboom2022equivariant}和离散数据建模\cite{gu2022vector}等领域。

除了强大的生成能力，扩散模型还具备内禀的去噪能力——通过迭代的噪声估计与去除过程，能够从含噪观测中有效恢复原始数据分布\cite{10419041}。这一特性为将扩散模型的去噪机制迁移至判别式任务提供了理论基础。

\section{判别式表征学习}

与生成式模型关注数据分布$P(\mathbf{x})$或联合分布$P(\mathbf{x}, y)$的建模不同，判别式模型专注于学习条件概率分布$P(y|\mathbf{x})$，其中$\mathbf{x}$表示输入数据，$y$表示任务特定的输出或标签。根据任务类型的不同，判别式模型涵盖了广泛的应用：分类模型\cite{agarap2017architecture,albawi2017understanding,deng2012mnist}将输入映射为离散的类别标签；检索模型\cite{lin2015deep,weyand2020google}学习特征空间使语义相似的样本在特征空间中彼此靠近；检测模型\cite{law2018cornernet,ren2015faster}则通过预测空间位置和边界框来定位和分类目标。

表征学习是各类判别式任务的核心基础，其目标是学习数据的有效表示以便于下游任务的信息提取\cite{6472238}。在视觉表征学习领域，骨干网络的架构设计经历了数次重要的范式变革。

卷积神经网络（CNN）是深度视觉表征学习的奠基架构。LeNet\cite{lecun1998gradient}首次将卷积运算应用于图像识别，奠定了CNN的基本结构。AlexNet\cite{krizhevsky2012imagenet}在2012年ImageNet\cite{deng2009imagenet}大规模视觉识别挑战赛上取得了突破性成绩，标志着深度学习时代的开启。此后，一系列精心设计的CNN架构不断推动性能的提升。ResNet\cite{he2016deep}通过引入残差连接解决了深层网络的梯度消失问题，使训练数百层深度的网络成为可能。RegNet\cite{radosavovic2020designing}通过系统性的网络设计空间搜索，提出了一组参数化的网络设计准则。EfficientNet\cite{tan2019efficientnet}通过复合缩放策略同时优化网络的深度、宽度和分辨率，在多个基准上实现了精度与效率的良好平衡。值得注意的是，在Transformer架构主导视觉领域之后，Liu等人\cite{liu2022convnet}提出了ConvNeXt，通过系统性地将现代设计元素（如大卷积核、层归一化、GELU激活等）引入标准卷积网络，在ImageNet分类上达到87.8\%的top-1精度，并在COCO检测和ADE20K分割任务上超越了Swin Transformer，证明了纯卷积架构仍然具有与Transformer相当的竞争力。

自注意力机制的引入开启了视觉表征学习的新范式。Transformer架构最初由Vaswani等人\cite{vaswani2017attention}为自然语言处理领域提出，其核心的多头自注意力机制能够捕捉输入序列中任意位置之间的长程依赖关系。Dosovitskiy等人\cite{dosovitskiy2020image}将Transformer直接应用于图像分类，提出了Vision Transformer（ViT），将图像划分为一系列固定大小的图像块（Patch），并将每个图像块视为一个词元（Token）进行处理。ViT在大规模数据集上取得了与CNN相当甚至更优的性能，证明了Transformer架构在视觉任务中的强大潜力。DeiT\cite{touvron2021training}通过精心设计的知识蒸馏策略和数据增强方案，使ViT在中等规模数据集上也能高效训练。Swin Transformer\cite{liu2021swin}引入了层级结构和移位窗口注意力机制，在保持线性计算复杂度的同时实现了跨窗口信息交互，成为多种视觉任务的通用骨干网络。

在自监督表征学习方面，近年来涌现了一系列重要工作。He等人\cite{he2022masked}提出了掩码自编码器（MAE），通过随机掩盖75\%的图像块并重建缺失像素来进行自监督预训练。MAE采用非对称的编码器-解码器架构，其中编码器仅处理可见的图像块子集，从而实现了高效的大规模模型训练。在ImageNet-1K上，ViT-Huge模型通过MAE预训练达到了87.8\%的top-1精度，同时在下游任务上展现出优异的迁移性能。Oquab等人\cite{oquab2023dinov2}提出的DINOv2进一步推进了自监督视觉基础模型的发展——通过在经过自动筛选的大规模多源图像数据上训练ViT模型（参数量达10亿），并结合自蒸馏策略将知识迁移至一系列小型模型，DINOv2能够产生在不同图像分布和任务间具有强泛化能力的通用视觉特征，在图像级和像素级的多个基准上超越了此前的最优方法。

最近，基于状态空间模型的VMamba\cite{liu2024vmamba}引入了选择性扫描机制，能够以线性复杂度建模长序列依赖关系，在图像分类、目标检测和语义分割等任务上展现出与Transformer架构相当的性能，同时具备更高的计算效率。

在具体的判别式任务中，行人重识别（Person Re-Identification, ReID）\cite{ye2021deep,bedagkar2014survey}是一个对特征判别性要求极高的代表性任务。ReID旨在跨不相交摄像头视角匹配同一行人的图像，面临姿态变化、光照差异和遮挡等严峻挑战。ReID方法的发展主要经历了两个阶段：早期方法依赖手工设计的描述子\cite{liao2015person,ma2014covariance,yang2014salient}结合度量学习\cite{koestinger2012large,liao2015efficient,zheng2013reidentification}；现代方法则以基于深度学习的端到端框架为主\cite{cho2022part,fu2021unsupervised,ni2023part,wang2023self,wang2019color,wang2020high}。当前的ReID系统普遍采用CNN\cite{lecun1998gradient}或Transformer\cite{dosovitskiy2020image,vaswani2017attention}作为骨干网络建模复杂的空间关系与层级化视觉表示，并通过注意力机制\cite{dosovitskiy2020image,vaswani2017attention}、时空推理\cite{li2022autost,li2020dyngcn}和域适应方法\cite{chen2018domain}进一步增强模型在实际部署场景中的鲁棒性和泛化能力。此外，RepVGG\cite{ding2021repvgg}提出的结构重参数化思想——在训练阶段使用多分支拓扑以增强表达能力，在推理阶段将多分支融合为单一卷积以保持高效——为本文的参数融合策略提供了重要的技术启发。

\section{扩散模型在判别任务中的应用}

尽管扩散模型在生成任务中取得了巨大成功，其在判别式任务中的探索仍处于起步阶段。近年来，一些研究开始尝试将扩散机制引入特定的判别式任务框架中。

在目标检测领域，DiffusionDet\cite{chen2023diffusiondet}首次将目标检测建模为从噪声边界框到精确目标框的去噪扩散过程。该方法在训练阶段对真实目标框施加前向扩散以生成噪声框作为训练样本，在推理阶段则从随机生成的噪声框出发，通过迭代去噪逐步恢复精确的目标框坐标。DiffusionDet在与Faster R-CNN\cite{ren2015faster}等成熟检测器的比较中取得了具有竞争力的表现，展示了扩散范式在检测任务中的可行性。

在图像分割领域，DiffSeg\cite{tian2023diffuse}利用预训练的稳定扩散模型实现了无监督零样本图像分割。该方法提出了一种迭代注意力合并机制，通过聚合基于KL散度的自注意力图来生成有效的分割掩码，无需额外训练数据或文本监督信号即可获得良好的分割效果。此外，Baranchuk等人\cite{baranchuk2022label}从另一个角度验证了扩散模型在分割任务中的潜力——他们发现，预训练扩散模型在反向去噪过程中的中间层激活能够有效捕捉输入图像的语义信息，构成优质的像素级表征。基于此发现，他们提出了一种简洁的少样本分割方法，仅需少量标注图像即可在多个数据集上显著超越已有的替代方案。这一工作从实证角度表明，扩散模型内部的去噪特征天然蕴含了丰富的判别性语义信息。

在零样本分类领域，Li等人\cite{li2023diffusionclassifier}提出了扩散分类器（Diffusion Classifier）的概念，证明了大规模文本到图像扩散模型（如Stable Diffusion）的条件密度估计可以直接用于执行零样本分类，无需任何额外训练。该方法通过比较不同类别文本提示下扩散模型对输入图像的条件似然来进行分类决策，在多种基准上取得了优异的结果，且在多模态组合推理能力上显著优于现有的判别式方法。在零样本学习领域，CLIP-Diffusion\cite{li2023clip}探索了将预训练的视觉-语言模型CLIP与扩散机制相结合的方案，利用扩散过程增强视觉特征与文本描述之间的对齐，在零样本分类和检索任务中取得了改进。

然而，上述方法均存在共同的局限性。首先，这些方法针对特定任务进行定制化设计，依赖于特定的数据结构——例如DiffusionDet操作的是噪声边界框空间，DiffSeg操作的是噪声初始化的分割图——因此难以直接迁移至其他判别式任务，缺乏通用性。其次，多步递归的去噪推理过程引入了显著的额外计算开销，在实际部署中面临效率瓶颈。这些局限性促使研究者探索更为通用且高效的扩散式特征增强范式。

\section{高效扩散采样方法}

扩散模型面临的核心挑战之一在于推理阶段的高计算成本，因为传统的去噪过程通常需要数百乃至数千步迭代采样才能生成高质量样本。为缓解这一效率瓶颈，近年来涌现了一系列加速采样技术。

DDPM\cite{ho2020denoising}建立了基本的迭代去噪框架，虽然能够生成高质量样本，但推理时间较长。为提升效率，Song等人\cite{song2020denoising}提出了去噪扩散隐式模型（Denoising Diffusion Implicit Model, DDIM），通过引入确定性的非马尔可夫采样过程重新参数化了反向去噪过程。DDIM在保持生成质量的前提下，将所需的采样步数从原始DDPM的数千步大幅减少至数十步，显著提升了推理效率。

在DDIM的基础上，渐进式蒸馏（Progressive Distillation）\cite{salimans2022progressive}提出了一种基于知识蒸馏的采样加速策略。该方法通过训练一个学生模型来匹配教师模型执行两步连续去噪的输出，从而在每一轮蒸馏中将所需采样步数减半。通过多轮迭代蒸馏，可以将长扩散轨迹有效压缩为更少的采样步数，且不显著损失生成质量。这种蒸馏式加速策略为本文第四章中设计的多步去噪蒸馏机制提供了直接的技术启发。

DPM-Solver\cite{lu2022dpm}从另一个角度解决了高效采样问题。该方法将扩散过程建模为常微分方程（ODE），并采用高阶数值求解器来高效近似去噪轨迹。DPM-Solver仅需10至20步采样即可生成高保真度的样本，在采样效率上实现了进一步突破。

Song等人\cite{song2023consistency}提出了一致性模型（Consistency Models），从根本上改变了扩散采样的范式。不同于传统方法通过多步迭代逐步去噪，一致性模型通过学习概率流ODE轨迹上任意点到原始数据的直接映射，实现了单步或少步高质量采样。一致性模型既可以通过蒸馏预训练的扩散模型获得，也可以作为独立的生成模型从头训练，在CIFAR-10上取得了单步生成3.55的FID，在ImageNet 64$\times$64上取得了6.20的FID，均刷新了当时的单步生成记录。Luo等人\cite{luo2023latent}进一步将一致性模型扩展至隐空间，提出了隐空间一致性模型（Latent Consistency Models, LCM），能够在任意预训练的隐空间扩散模型（如Stable Diffusion）上进行高效蒸馏，仅需2至4步即可生成高质量的768$\times$768分辨率图像，整个训练过程仅需32块A100 GPU小时。

Liu等人\cite{liu2022rectified}提出了矫正流（Rectified Flow），通过学习连接噪声分布与数据分布之间的直线路径，构建了一种简洁高效的ODE生成模型。矫正流的核心思想是在两个分布的样本之间学习尽可能直的传输路径——直线路径不仅是两点间的最短路径，而且可以在不进行时间离散化的情况下精确模拟，因此产生了计算效率极高的模型。通过递归地应用矫正操作，可以获得路径越来越直的流模型，使得粗粒度的时间离散化（甚至单步Euler离散化）也能产生高质量的生成结果。

上述加速采样技术的发展反映了扩散模型领域在去噪精度与计算效率之间寻求平衡的持续努力。这些工作，特别是基于蒸馏的加速方法和一致性学习范式，为本文在有限计算预算下设计高效的特征去噪增强机制提供了重要的方法论参考。

\section{知识蒸馏相关方法}

知识蒸馏的概念最早由Hinton等人\cite{hinton2015distilling}系统性地提出，其核心思想是将一个参数量较大、表达能力较强的教师模型所学到的``暗知识"（Dark Knowledge）迁移给一个参数量较小的学生模型。具体而言，教师模型输出的软化概率分布（通过温度系数$\tau$调节softmax函数）编码了类别之间的相似性关系，这些信息比硬标签更加丰富，能够为学生模型提供更加细粒度的指导。知识蒸馏已在模型压缩、增量学习和跨模态迁移学习等多种场景中得到了广泛应用\cite{8395024}。

在蒸馏方法的发展过程中，Romero等人\cite{romero2015fitnets}提出了FitNets，将知识蒸馏从输出层扩展到中间层特征表示。FitNets利用教师网络的中间隐层表征作为``提示"（Hints）来指导学生网络对应层的训练，并引入适配层来处理教师与学生中间层之间的维度差异。这种基于特征的蒸馏方法使得学生网络可以更深更窄，在CIFAR-10上仅用教师约十分之一的参数量即可超越教师网络的性能，开创了特征级知识蒸馏的研究方向。

Tian等人\cite{tian2020contrastive}从表征学习的角度重新审视了知识蒸馏，提出了对比表征蒸馏（Contrastive Representation Distillation, CRD）。该方法指出传统的KL散度蒸馏损失只关注教师与学生输出分布的匹配，而忽略了教师表征中蕴含的结构化知识。CRD通过对比学习的方式训练学生模型，使其能够捕获教师表征中更丰富的结构信息——学生特征应当与对应的教师特征在表征空间中彼此靠近，同时远离其他样本的教师特征。CRD在多种知识迁移场景中均超越了传统蒸馏方法。

在视觉表征学习领域，知识蒸馏被广泛用于训练高效的小型模型。例如，DeiT\cite{touvron2021training}通过引入蒸馏令牌（Distillation Token），利用CNN教师模型指导ViT学生模型的训练，使得Vision Transformer在没有大规模预训练数据的条件下也能取得优异性能。在行人重识别任务中，自蒸馏和互蒸馏策略也被用于增强特征的判别性\cite{luo2021self}。

在扩散模型的语境下，知识蒸馏同样发挥着重要作用。如前节所述，渐进式蒸馏\cite{salimans2022progressive}通过教师-学生框架将多步采样过程压缩为更少的步数，实现了扩散模型推理的显著加速。一致性模型\cite{song2023consistency}亦可通过蒸馏预训练扩散模型来获得单步生成能力。这些工作为本文的研究提供了重要启发：当去噪模块为了实现推理阶段的参数融合而被设计为轻量化结构时，其有限的参数量会制约对复杂高维噪声分布的建模能力。此时，可以借助一个表达能力更强的教师网络在训练阶段提供更精确的噪声分布估计，通过蒸馏将教师的去噪知识传递给轻量化的学生去噪模块，从而在不增加推理开销的前提下提升去噪质量。

\section{本章小结}

本章从五个方面综述了与本文研究密切相关的理论基础与技术方法。在生成式扩散模型方面，回顾了从GAN、VAE到扩散模型的发展历程，详细阐述了扩散模型的前向扩散与反向去噪数学框架，以及DDPM、改进型DDPM、基于分数的SDE等里程碑式工作，并介绍了DiT和流匹配等前沿的架构与范式创新。在判别式表征学习方面，梳理了从CNN到Vision Transformer再到状态空间模型的骨干网络演进脉络，涵盖了ConvNeXt等现代卷积架构的回归、MAE和DINOv2等自监督学习方法的突破，以及行人重识别等典型任务中的方法发展。在扩散模型的判别式应用方面，分析了DiffusionDet、DiffSeg、扩散分类器和扩散特征语义分割等代表性工作的技术方案及其在通用性和效率方面的局限。在高效扩散采样方面，介绍了DDIM、渐进式蒸馏、DPM-Solver、一致性模型、隐空间一致性模型和矫正流等加速技术，为本文的多步去噪融合策略提供方法论参考。在知识蒸馏方面，从Hinton等人的开创性工作出发，阐述了FitNets特征级蒸馏和CRD对比表征蒸馏等重要方法，以及知识蒸馏在视觉表征学习和扩散模型加速中的应用，为本文的蒸馏增强去噪模块奠定理论基础。上述综述为后续第三章和第四章的方法设计提供了全面的技术参照与理论支撑。